\chapter{Gasto federalizado}

\titular{Participaciones y aportaciones suman \textbf{2\,497\,938,8 millones de pesos},
\textbf{6,45\,\% del PIB} y casi la cuarta parte del gasto neto total. Las participaciones
crecen \textbf{3,7\,\% real} y las aportaciones \textbf{1,5}. \textbf{Lo que se transfiere sin
condiciones crece a más del doble de la tasa de lo etiquetado} ---tasas, no montos: el ramo
etiquetado también aumenta---, y esa diferencia es la política federalista del paquete.}

\quecambio{El Ramo 28 pasa de 1\,340\,210,9 a 1\,456\,045,9 millones y el Ramo 33 de
979\,951,8 a 1\,041\,892,9. Dentro del 33, el fondo educativo crece 0,1\,\% real y el de salud
cae 0,6. Y todo el ramo cambió de unidad responsable, de la 420 a la 411.}

\porqueimporta{Las participaciones son ingreso propio de las entidades y las aportaciones son
gasto federal etiquetado que ellas ejecutan. \textbf{Que las primeras crezcan 3,7\,\% y las
segundas 1,5 significa que los recursos de libre disposición crecen más rápido que los
condicionados}, no que los condicionados bajen: las aportaciones para educación y salud
---los dos destinos mayores del ramo etiquetado--- aumentan en términos reales, a menor tasa.}

\section{Declaración de perímetro}

\marginal{dos vías, dos objetos}
Este capítulo mide dos ramos generales y no los suma con nada más: el \textbf{28},
participaciones a entidades federativas y municipios, y el \textbf{33}, aportaciones federales.
Quedan fuera los convenios de reasignación, que se pactan durante el ejercicio y no están en el
proyecto, y las transferencias que otros ramos hacen a los estados sin ser aportaciones.

\textbf{Las dos vías no son el mismo objeto}: las participaciones son recursos que la
Constitución reconoce como propios de las entidades y las aportaciones son gasto federal
etiquetado que ellas ejecutan. Se presentan juntas porque el lector quiere saber cuánto baja a
los estados, y se etiquetan siempre porque no se pueden discutir igual.

\section{Las dos vías}

\begin{list}{}{\setlength{\leftmargin}{1.4em}\setlength{\itemsep}{2pt}}
\item[$\bullet$] \textbf{Participaciones (Ramo 28):} 1\,456\,045,9 millones, \textbf{3,7\,\%
real}, 3,76\,\% del PIB. Dentro, el Fondo General de Participaciones pasa de 980\,334,4 a
1\,069\,854,4 millones.
\item[$\bullet$] \textbf{Aportaciones (Ramo 33):} 1\,041\,892,9 millones, \textbf{1,5\,\%
real}, 2,69\,\% del PIB. El fondo de nómina educativa crece 0,1\,\% real y el de servicios de
salud cae 0,6.
\end{list}

\textbf{Cinco décimas de punto del PIB separan hoy a las dos vías} y la brecha se ensancha:
las participaciones dependen de la recaudación federal participable, que crece con la meta
tributaria de 15,1\,\% del PIB; las aportaciones dependen de decisiones de gasto que este año
las mantuvieron planas.

\subsection{El ramo entero cambió de unidad responsable}

Todo el Ramo 33 pasó de la unidad 420 a la \textbf{411}. Son 1\,041\,892,9 millones de pesos
cambiando de unidad responsable sin cambiar de destino. Los ramos 19 y 30 hicieron el mismo
movimiento. \textbf{Cualquier serie construida por unidad responsable se rompe este año}, y
ninguna decisión de gasto está detrás.

\section{Por entidad: lo que el proyecto sí identifica y lo que no}

\marginal{la columna territorial no es un reparto}
El proyecto identifica una entidad federativa para parte de su gasto, y esa columna
\textbf{no es un reparto del gasto federalizado}. Dos cifras lo demuestran: a la Ciudad de
México se le atribuyen \textbf{4\,907\,088,4 millones de pesos} ---casi la mitad del
presupuesto--- y hay \textbf{1\,731\,677,0 millones} marcados como no distribuibles
geográficamente.

La primera cifra es el gasto de administración central, que se registra donde está la
dependencia. \textbf{Dividirla entre los habitantes de la capital produciría una cifra per
cápita absurda}, y es exactamente el tipo de error que este documento existe para no cometer.

\textbf{Por eso no se publica gasto federalizado por habitante y por entidad.} El reparto
efectivo de los fondos del Ramo 33 está en el Anexo 22 del decreto, que sí lo trae fondo por
fondo, y su explotación queda fuera del alcance de esta corrida.

\section{Implicaciones}

El paquete transfiere más a los estados y condiciona menos. Visto desde el federalismo fiscal
es una buena noticia; visto desde los servicios que las aportaciones financian, no lo es:
\textbf{la nómina educativa se queda plana en términos reales y el fondo de salud sigue en el
nivel bajo al que cayó hace dos ejercicios.}

Y hay un punto que este capítulo no puede cerrar y conviene enunciar. \textbf{Las
participaciones crecen porque se supone que la recaudación crecerá 5,2\,\% real}, y la mitad
de ese aumento viene de comercio exterior y de accesorios de impuestos. Si esa meta no se
cumple, las entidades reciben menos por la vía que este año más creció, y lo reciben durante
el ejercicio, sin margen para ajustar.

\begin{ficha}{gasto federalizado}
\fila{Perímetro}{Ramos generales 28 y 33, presentados por separado y nunca sumados sin
etiqueta. Convenios de reasignación y transferencias de otros ramos quedan fuera.}
\fila{Línea de comparación}{\textbf{Línea P:} proyecto 2026 contra proyecto 2025.}
\fila{Deflactor}{Del PIB, 4,8\,\%, CGPE 2026.}
\fila{Año de los pesos}{Corrientes; diferencias reales en pesos de 2026.}
\fila{PIB y añada}{38\,715,9 miles de millones, CGPE 2026.}
\fila{Denominador demográfico}{\textbf{No se publica ninguna cifra per cápita en este
capítulo}, por la razón que se explica en el cuerpo: la columna territorial del proyecto no es
un reparto.}
\fila{Fuentes}{Analíticos del proyecto por programa, 2025 y 2026; Anexo 1 del proyecto de
decreto; base del proyecto por entidad federativa.}
\fila{Qué no puede afirmarse}{El gasto federalizado por habitante y por entidad. El reparto de
los fondos del Ramo 33 entidad por entidad, que exige explotar el Anexo 22 y no se hizo. Y los
convenios de reasignación, que no están en el proyecto por definición.}
\end{ficha}
\clearpage
